\documentclass[11pt]{article}
\usepackage[a4paper,margin=25mm]{geometry}
\usepackage{amsmath,amssymb,amsthm}
\usepackage[hidelinks]{hyperref}
\newtheorem{theorem}{Theorem}
\newtheorem{lemma}{Lemma}
\theoremstyle{remark}
\newtheorem{remark}{Remark}
\newcommand{\F}{\mathbb F}
\title{Matrices with Squarefree Characteristic Polynomial:\\A Disproof of Conjecture 00000001854}
\author{GitHub: gaochengzhecpu}
\date{}
\begin{document}
\maketitle
\begin{abstract}
The source asserts that the number of $n\times n$ matrices over $\F_q$
with squarefree characteristic polynomial is exactly
$q^{n^2}\prod(1-q^{-i^2})$. For the product over $1\le i\le n$ we show
that the identity fails for every finite field and every $n\ge1$: for
$n=1$ the count is $q$ and the formula gives $q-1$, and for $n\ge2$ the
formula is not an integer. At $n=1$ the identity also fails for every
non-empty truncation of the product and for the infinite product, and
the product over $1\le i\le n-1$ fails at $n=5$. These
statements are verified in Lean with Mathlib's matrices, characteristic
polynomials and squarefree elements.
\end{abstract}

\paragraph{Notation and reading of the source.}
For a finite field $F$ with $q$ elements and $n\ge0$ let
\[
 N_n(F)=\#\{M\in \mathrm{Mat}_n(F): \chi_M \text{ is squarefree}\},
\]
where $\chi_M=\det(XI-M)\in F[X]$ and a polynomial is squarefree if no
square of a non-unit divides it. Both language versions of the source
state
\[
 N_n(\F_q)=q^{n^2}\prod\bigl(1-q^{-i^2}\bigr)
\]
as an exact closed form, without giving the range of the index $i$. The
shape of the formula imitates $|\mathrm{GL}_n(\F_q)|=q^{n^2}\prod_{i=1}^n
(1-q^{-i})$, so the natural reading is the product over $1\le i\le n$
(Reading~A). The other reading we treat is the infinite product over
$i\ge1$ (Reading~B). The claim is read as an identity for all prime
powers $q$ and all $n\ge1$.

\begin{theorem}\label{thm:one}
For every finite field $F$ with $q$ elements, $N_1(F)=q$.
\end{theorem}
\begin{proof}
The characteristic polynomial of the $1\times1$ matrix $(a)$ is $X-a$.
A polynomial of degree one over a field is irreducible, and an
irreducible element is squarefree: if $d^2$ divides $X-a$ then $d$
divides an irreducible element twice, which forces $d$ to be a unit.
Hence all $q$ matrices are counted.
\end{proof}

\begin{theorem}[Reading A]\label{thm:A}
Let $F$ be a finite field with $q$ elements and let $n\ge1$. Then
\[
 N_n(F)\ne q^{n^2}\prod_{i=1}^{n}\bigl(1-q^{-i^2}\bigr).
\]
For $n=1$ the right-hand side is $q-1$, and for $n\ge2$ it is not an
integer.
\end{theorem}
\begin{proof}
For $n=1$ the right-hand side is $q(1-q^{-1})=q-1$, while $N_1(F)=q$ by
Theorem~\ref{thm:one}.

Let $n\ge2$ and put $S=\sum_{i=1}^n i^2$. Suppose that the right-hand
side equals an integer $N$. Multiplying by $q^{S}=\prod_{i=1}^n q^{i^2}$
gives the identity of integers
\[
 N\,q^{S}=q^{n^2}\prod_{i=1}^{n}\bigl(q^{i^2}-1\bigr).
\]
Since $n\ge2$, the sum $S$ contains the two different terms $1^2$ and
$n^2$, so $S\ge n^2+1$. Cancelling $q^{n^2}$ shows that $q$ divides
$P=\prod_{i=1}^n(q^{i^2}-1)$. But every factor is coprime to $q$,
because $q^{i^2-1}\cdot q-(q^{i^2}-1)=1$ for $i\ge1$; hence $P$ is
coprime to $q$. A common divisor $q$ of $q$ and $P$ must then be a
unit, contradicting $q\ge2$.
\end{proof}

\begin{theorem}[Other finite ranges]\label{thm:C}
Let $q\ge2$ and let $n,m\ge1$ with $\sum_{i=1}^{m}i^2>n^2$. Then
$q^{n^2}\prod_{i=1}^{m}(1-q^{-i^2})$ is not an integer. In particular,
with the product over $1\le i\le n-1$ the identity fails at $n=5$ for
every finite field, because $1+4+9+16=30>25$.
\end{theorem}
\begin{proof}
The proof of Theorem~\ref{thm:A} for $n\ge2$ used only $S>n^2$, and it
applies with $S=\sum_{i=1}^{m}i^2$ and the product over $1\le i\le m$.
\end{proof}

\begin{theorem}[Truncations and Reading B at $n=1$]\label{thm:B}
Let $F$ be a finite field with $q$ elements.
\begin{enumerate}
\item For every $m\ge1$,
 $q\prod_{i=1}^{m}(1-q^{-i^2})\le q-1<q=N_1(F)$.
\item If the partial products $\prod_{i=1}^{m}(1-q^{-i^2})$ converge to
 $L$ as $m\to\infty$, then $qL\le q-1<N_1(F)$.
\end{enumerate}
\end{theorem}
\begin{proof}
Every factor $1-q^{-i^2}$ lies in $[0,1]$, so the partial products are
non-negative and non-increasing in $m$. For $m\ge1$ they are therefore
at most the first factor $1-q^{-1}$, which gives (1). The limit of a
sequence that is eventually at most $1-q^{-1}$ is at most $1-q^{-1}$,
which gives (2).
\end{proof}

Thus the asserted identity is false in Reading~A for every $q$ and
every $n\ge1$, and in Reading~B at $n=1$ for every $q$.

\begin{remark}[The true values for small $n$; not formalised]
For $n=2$ one has $N_2(\F_q)=q^4-q^3$. Indeed, a monic quadratic is not
squarefree exactly when it equals $(X-a)^2$ for some $a\in\F_q$, and
$a$ is then unique. A matrix $M$ has $\chi_M=(X-a)^2$ exactly when
$M-aI=\bigl(\begin{smallmatrix}x&y\\z&-x\end{smallmatrix}\bigr)$ with
$x^2+yz=0$. This equation has $q^2$ solutions: $q(q-1)$ with $y\ne0$
(then $z=-x^2/y$) and $q$ with $y=0$ (then $x=0$). So $q\cdot q^2$
matrices are excluded. Reading~A gives $(q^4-q^3)(1-q^{-4})$ instead.
The supplementary script \texttt{verify.py} enumerates all
matrices for small parameters and finds $N_1=q$ and $N_2=q^4-q^3$ for
$q\in\{2,3,5\}$, $N_3(\F_2)=160$ and $N_3(\F_3)=11178$. In particular a
product over $1\le i\le n-1$, which happens to agree with $N_n$ for
$n\le2$, gives $240\ne160$ at $(n,q)=(3,2)$. This remark is supported by
the paper argument and the enumeration only; it is not used in the
theorems above.
\end{remark}

\section*{The earlier submission and its error}
Pull request \#23 (orionsheep) treated this conjecture with the same
counterexample at $n=1$ and the same two readings. It was merged on
2026-10-01 and removed by the organisers in commit \texttt{4504549}
the same day; the stated reason was that its Lean theorems carried none
of the mathematical objects. The mathematics of that submission was
correct. The error was in its Lean file, which used Lean core only. The
squarefree predicate was defined as the constant function
\begin{center}\small\texttt{def sqfree1b : Fin 2 -> Bool := fun \_ => true},\end{center}
the count was the length of a two-element list filtered by this
predicate, and the decisive statements were closed arithmetic facts
such as \texttt{(2 : Nat) $\neq$ 1}. No matrix, characteristic
polynomial, squarefreeness, finite field or quantifier over $q$ and $n$
occurred. The present submission defines the count with Mathlib's
\texttt{Matrix.charpoly} and \texttt{Squarefree}, proves $N_1(F)=q$ for
every finite field, and proves the failure of the identity for all
finite fields and all $n\ge1$.

\section*{Formal verification and scope}
The Lean definition
\begin{center}\small\texttt{sqfreeCount F n}\end{center}
is, for an arbitrary field \texttt{F}, the cardinality of the subtype of
\begin{center}\small\texttt{Matrix (Fin n) (Fin n) F}\end{center}
on which \texttt{Squarefree M.charpoly} holds. The right-hand side is
\texttt{formula q n m}, the rational number
$q^{n^2}\prod_{i=1}^{m}(1-1/q^{i^2})$, and
\texttt{partialProduct q m} is the real partial product.

The theorems
\begin{center}\small
\texttt{squarefree\_charpoly\_one},\qquad \texttt{sqfreeCount\_one}
\end{center}
are Theorem~\ref{thm:one} for every finite field. The theorem
\begin{center}\small\texttt{formula\_not\_natural}\end{center}
states that for $q\ge2$ and $n\ge2$ no natural number equals
\texttt{formula q n n}. The theorem \texttt{formula\_fails} is
Theorem~\ref{thm:A}: for every type \texttt{F} with \texttt{Field} and
\texttt{Fintype} instances and every $n\ge1$,
\begin{center}\small
\texttt{(sqfreeCount F n : Rat) $\neq$ formula (Fintype.card F) n n}.
\end{center}
Theorem~\ref{thm:C} is
\begin{center}\small
\texttt{formula\_not\_natural\_of\_sum},\qquad \texttt{shifted\_formula\_fails},
\end{center}
the second being the case $n=5$, $m=4$ for every finite field.
Theorem~\ref{thm:B}(1) is
\begin{center}\small\texttt{truncated\_formula\_fails\_at\_one},\end{center}
and Theorem~\ref{thm:B}(2) is
\begin{center}\small\texttt{infinite\_product\_fails\_at\_one},\end{center}
with convergence expressed by \texttt{Filter.Tendsto}. The propositions
\begin{center}\small
\texttt{ClaimedClosedForm},\qquad \texttt{ClaimedClosedFormInfinite}
\end{center}
are the universally quantified identities of Readings~A and~B. Their
negations are the final theorems
\begin{center}\small
\texttt{conjecture\_false},\qquad \texttt{conjecture\_false\_infinite}.
\end{center}

The negations need one concrete finite field. The Mathlib module that
declares the field structure on \texttt{ZMod p} is not among the
modules imported here, so the file declares the instance
\texttt{instFieldZModTwo} from Mathlib's commutative ring structure and
inversion on \texttt{ZMod 2}, proving $a\,a^{-1}=1$ for $a\ne0$. The
general theorems do not depend on this instance. No
\texttt{native\_decide}, \texttt{sorry} or custom axiom occurs.

Scope. The source does not fix the index range of the product. The
formal results cover the product over $1\le i\le n$ for all $n\ge1$,
and at $n=1$ every truncation $1\le i\le m$ with $m\ge1$ and the
infinite product. Reading~B is refuted only at $n=1$, which suffices
for a claim about all $n$. The product over $1\le i\le n-1$, which is
empty at $n=1$ and agrees with $N_n$ for $n\le2$, is refuted formally
at $n=5$ by Theorem~\ref{thm:C}. More generally a range
$1\le i\le m(n)$ is refuted as soon as $m(1)\ge1$, or
$\sum_{i\le m(n)}i^2>n^2$ for some $n$. Index ranges that do not start
at $i=1$ are not covered. No claim is made about the asymptotic
proportion of such matrices as $n$ or $q$ grows.

\section*{Source and reproduction}
The exact bilingual source is included byte-for-byte as
\texttt{SOURCE.md}. Its repository path is
\begin{center}\small\texttt{conjectures/00000001854.md}.\end{center}
The project pins Lean 4.19.0 and Mathlib commit
\begin{center}\small\texttt{c44e0c8ee63ca166450922a373c7409c5d26b00b}.\end{center}
From \texttt{lean/}, run \texttt{lake build}, then
\begin{center}\small\texttt{lake env lean -DwarningAsError=true Main.lean}.\end{center}
The included public-Git manifest locks all dependency revisions.
The \texttt{verification/} directory records the fresh-build logs,
theorem-axiom reports, artifact hashes and the self-review. The script
\texttt{verify.py} is supplementary and supports only the remark.
\end{document}
